\documentclass[11pt,a4paper]{article}
\usepackage[T1]{fontenc}
\usepackage{lmodern,amsmath,amssymb,amsthm}
\usepackage[margin=27mm]{geometry}
\usepackage[colorlinks=true,linkcolor=blue,citecolor=blue,urlcolor=blue]{hyperref}
\usepackage{microtype}
\newtheorem{theorem}{Theorem}[section]
\newtheorem{proposition}[theorem]{Proposition}
\newtheorem{lemma}[theorem]{Lemma}
\newtheorem{corollary}[theorem]{Corollary}
\theoremstyle{definition}\newtheorem{definition}[theorem]{Definition}
\newtheorem{remark}[theorem]{Remark}
\title{An obstruction to one-parameter presentations of an elliptic point}
\author{}\date{7 November 2026}
\begin{document}\maketitle
\begin{abstract}
For $K=\mathbb Q(\sqrt2)$ and
$E:y^2=x^3-(11-6\sqrt2)x$, every algebraic curve over
$\mathbb Q$ mapping to $\operatorname{Res}_{K/\mathbb Q}E$
has finite rational image. Consequently no finite collection of
one-parameter algebraic presentations with finite fibers can represent
an infinite subset of $E(K)$. Two rational parameters suffice,
with one finite quartic choice and finitely many charts. The obstruction
is unconditional and concerns presentations of the elliptic point
itself; it is not a lower bound for arbitrary Diophantine definitions
of the integers.
\end{abstract}

\section{The elliptic curve and its conjugate}

Put
\[
 K=\mathbb Q(\sqrt2),\quad d=3-\sqrt2,\quad D=d^2=11-6\sqrt2,
 \qquad E:\ y^2=x^3-Dx.
\]
Let $\sigma$ be the nontrivial automorphism of $K/\mathbb Q$,
and write $E^\sigma:y^2=x^3-\sigma(D)x$.

\begin{proposition}\label{prop:nonisogenous}
The elliptic curves $E$ and $E^\sigma$ are not isogenous over $K$.
More precisely,
\[
                  \operatorname{Hom}_{K(i)}(E,E^\sigma)=0.
\]
\end{proposition}
\begin{proof}
The field $K(i)$ has a place with residue field $\mathbb F_{17}$
at which $\sqrt2$ reduces to $6$ and $i$ reduces to $4$.
At this place both curves have good reduction, and their equations are
\[
                 y^2=x^3-9x,\qquad y^2=x^3-13x.
\]
The discriminants are nonzero because they are $64\cdot9^3$ and
$64\cdot13^3$ modulo seventeen. The square residues modulo seventeen
are $0,1,2,4,8,9,13,15,16$. Enumeration of the seventeen abscissas gives
\[
\begin{array}{c|ccc|c}
D\bmod17 & x^3-Dx=0 & \text{nonzero square} & \text{nonsquare}
         & \#E_D(\mathbb F_{17})\\\hline
9&3&8&6&20\\
13&3&6&8&16.
\end{array}
\]
The last column includes the point at infinity. A nonzero homomorphism
between elliptic curves is an isogeny. An isogeny at a place of good
reduction gives an isogeny of the reductions, whose Frobenius
polynomials, and therefore their numbers of rational points, are equal.
The unequal counts exclude such a homomorphism over $K(i)$, hence
also over $K$. The accompanying script reproduces the counts by two
independent finite enumerations.
\end{proof}

The conjugate curves are nevertheless quadratic twists of one another.
Indeed $d\sigma(d)=7$, and
\[
          \delta=\frac{\sigma(d)}d=\frac7{d^2}
\]
represents the squareclass of seven. Over $K(\sqrt7)$ the map
$(x,y)\mapsto(\delta x,\delta\sqrt\delta\,y)$ is an isomorphism
$E\to E^\sigma$. Thus the nonisogeny assertion concerns the stated
fields; it is not a claim of geometric nonisogeny.

\section{Simplicity of the Weil restriction}

The Weil restriction $A=\operatorname{Res}_{K/\mathbb Q}E$ is the
abelian surface characterized by
\[
                   A(R)=E(R\otimes_{\mathbb Q}K)
\]
for every $\mathbb Q$-algebra $R$. In particular $A(\mathbb Q)=E(K)$,
and $A_K\simeq E\times E^\sigma$.

\begin{proposition}\label{prop:simple}
The abelian surface $A$ is simple over $\mathbb Q$.
\end{proposition}
\begin{proof}
Suppose $B\subset A$ were a nonzero proper abelian subvariety
over $\mathbb Q$. It would have dimension one. The two projections
of $B_K$ into $E\times E^\sigma$ are homomorphisms of elliptic
curves. If both were nonzero, both would be isogenies, making
$E$ and $E^\sigma$ isogenous over $K$, contrary to
Proposition~\ref{prop:nonisogenous}.
Exactly one projection is therefore nonzero. The other is zero,
so $B_K$ is one of the two factors. But the nontrivial Galois
automorphism interchanges those factors. Neither factor descends
to a $\mathbb Q$-subvariety of $A$, a contradiction.
\end{proof}

\begin{lemma}\label{lem:infinite-group}
The group $E(K)$ is infinite. In particular the point
\[
 P=\left(\frac{57-28\sqrt2}{8},
          \frac{-660+417\sqrt2}{32}\right)
\]
has infinite order.
\end{lemma}
\begin{proof}
Its coordinates satisfy the curve equation. The prime three is
unramified in $K$, with residue field $\mathbb F_9$, and the
curve has good reduction there. Normalize its valuation by $v(3)=1$.
The displayed point has $v(x(P))=0$ and $v(y(P))=1$, while
$3x(P)^2-D$ is a unit. The tangent slope therefore has valuation
$-1$, and the duplication formulas give
\[
 v(x(2P))=-2,\qquad v(y(2P))=-3,\qquad v(z(2P))=1,
 \qquad z=-x/y.
\]
The odd integral formal multiplication series
$[m](Z)=mZ+Z^3g_m(Z^2)$ now shows that multiplication by an integer
prime to three preserves a positive parameter valuation, while
multiplication by three raises it by one. The higher terms have
strictly larger valuation than the linear term. Consequently
$v(z([2n]P))=1+v_3(n)$ for $n\ne0$, and no nonzero multiple of
$2P$ is the identity. Therefore $P$ has infinite order.
\end{proof}

\section{Images of algebraic curves}

\begin{theorem}\label{thm:curve-images}
Let $F$ be a number field and let $B/F$ be an abelian variety
which has no elliptic abelian subvariety defined over $F$.
If $V$ is an algebraic set over $F$ of dimension at most one,
$U\subseteq V$ is open, and $f:U\to B$ is a morphism, then
$f(U(F))$ is finite.
\end{theorem}
\begin{proof}
There are finitely many irreducible components. A zero-dimensional
component contributes finitely many points. If an irreducible curve
is not geometrically integral, all its $F$-points lie in the finite
intersections of its distinct conjugate geometric components.
It therefore suffices to consider a geometrically integral curve.

Take its smooth projective normalization $C$. Outside finitely many
points the normalization is an isomorphism, so rational points lift.
The induced rational map $C\dashrightarrow B$ extends to a morphism
on $C$ because $B$ is proper. If $C$ has genus at least two,
Faltings's theorem makes $C(F)$ finite \cite{Faltings1983}.

If $C$ has genus zero and has an $F$-point, it is isomorphic to
$\mathbb P^1_F$. Every morphism from $\mathbb P^1$ to an abelian
variety is constant \cite[Chapter I, Proposition 3.9]{MilneAV}.
If $C(F)$ is empty there is nothing to prove.

Finally suppose $C$ has genus one and an $F$-point $c_0$.
Then $(C,c_0)$ is an elliptic curve. Subtracting $f(c_0)$ from
$f$ gives a homomorphism of abelian varieties
\cite[Chapter I, Corollary 1.2]{MilneAV}. If it were nonzero,
its image would be an elliptic abelian subvariety of $B$ defined
over $F$, contrary to the hypothesis. Hence $f$ is constant.
The finitely many points omitted in passing to the normalization
only add a finite set of image values.
\end{proof}

\begin{corollary}\label{cor:one-parameter}
Let $C_1,\ldots,C_r$ be algebraic sets over $\mathbb Q$, each
finite over a nonempty open subset of $\mathbb A^1_{\mathbb Q}$.
For any rational maps $f_j:C_j\dashrightarrow A$, the union
of the images of their rational points, wherever the maps are
defined, is finite. In particular it cannot contain an infinite
subset of $E(K)=A(\mathbb Q)$.
\end{corollary}
\begin{proof}
Every $C_j$ has dimension at most one. Apply
Theorem~\ref{thm:curve-images} and Proposition~\ref{prop:simple},
then take the finite union.
\end{proof}

The corollary applies in particular to a retained rational parameter
$t$ followed by any finite list of triangular monic equations
\[
 F_j(t,z_1,\ldots,z_j)=0,
 \qquad F_j\text{ monic of positive degree in }z_j,
\]
with rational functions of the resulting coordinates specifying
the point of $E(K)$. Nonzero denominator conditions restrict to
an open subset and do not change the conclusion. Arbitrarily high
degrees and finitely many alternative charts do not avoid the
finiteness obstruction.

\section{A sharp two-parameter presentation}

\begin{proposition}\label{prop:two-parameters}
The set $E(K)$ admits a finite collection of presentations using
two retained rational parameters and one finite algebraic root
of degree at most four.
\end{proposition}
\begin{proof}
Retain $a_0,a_1\in\mathbb Q$ and put $a=a_0+a_1\sqrt2$.
Write $a^3-Da=A_0+A_1\sqrt2$, where
\[
 \begin{aligned}
 A_0&=a_0^3+6a_0a_1^2-11a_0+12a_1,\\
 A_1&=3a_0^2a_1+2a_1^3-11a_1+6a_0.
 \end{aligned}
\]
Let $Y$ be a rational root of
\[
                 Y^4-4A_0Y^2+8A_1^2=0,
                 \qquad Y\ne0.
\]
Then
\[
                    (a,\;Y/2+(A_1/Y)\sqrt2)
\]
is a point of $E(K)$. Conversely it represents every point whose
ordinate has nonzero rational component.

Put $T_0=(0,0)$. If $Q=(a,b)$ is not a two-torsion point,
at least one of $Q,Q+T_0$ has a nonzero rational component in
its ordinate. To see this, suppose both ordinates were in
$\sqrt2\mathbb Q^\times$. The addition formula gives
$b(Q+T_0)/b(Q)=D/a^2\in\mathbb Q$. Thus
$a/d\in\mathbb Q\cup\sqrt2\mathbb Q$.
The equation $b^2=d^3(a/d)((a/d)^2-1)\in\mathbb Q^\times$
would then imply $d^3\in\mathbb Q\cup\sqrt2\mathbb Q$,
whereas $d^3=45-29\sqrt2$ has both components nonzero.
This is a contradiction.

Use the displayed quartic chart directly, or translate its output
by $T_0$. The remaining points are the identity and the three
nonzero two-torsion points, which are finite fixed alternatives.
This gives the asserted presentation.
\end{proof}

The one-parameter obstruction and the two-parameter construction
give an exact minimum for this point-presentation problem. They
do not give a lower bound of 29 for a Diophantine definition of
$\mathbb Z$ over $\mathbb Q$. In particular, the theorem does not
assume that several elliptic points must be encoded independently
in every possible defining polynomial.

The finite certificates are reproduced by the standalone script
\texttt{verify\_elliptic\_chart.py}, using only Python's standard
library. It checks the two point counts independently, the displayed
elliptic points and their valuations at three, the twist identities,
and examples of the quartic chart. The general finiteness and
sharpness assertions follow from the proofs above.

\begin{thebibliography}{9}
\bibitem{MilneAV} J. S. Milne, \emph{Abelian Varieties},
version 2.0, 2008, Chapters I and IV.
\url{https://www.jmilne.org/math/CourseNotes/AV.pdf}.
\bibitem{Faltings1983} G. Faltings,
Endlichkeitss\"atze f\"ur abelsche Variet\"aten \"uber Zahlk\"orpern,
\emph{Invent. Math.} \textbf{73} (1983), 349--366;
erratum, \textbf{75} (1984), 381.
\url{https://doi.org/10.1007/BF01388432}.
\end{thebibliography}
\end{document}
